\section{Discussion}

\subsection*{Note on citations}
Bracket-number citations from the donor extraction have been converted to \texttt{natbib} citations where possible.

\subsection{The Response Sector}

\textbf{(see PDF; response-sector symbol not recovered from DOCX extraction)} is the effective stress--energy required so that the metric (or potentials) inferred from galaxy phenomenology satisfies the Einstein equation with baryons. This definition is ontologically agnostic---it does not commit to a microscopic composition of spacetime. We avoid claiming spacetime ``is literally'' a fluid or solid; instead, we describe the framework as an EFT/constitutive closure candidate.

\subsection{Relation to MOND and TeVeS}

Modified Newtonian Dynamics (MOND) \citep{Milgrom1983} shares the empirical acceleration scale $a_0$ and produces flat rotation curves through a modified force law. The response-sector framework differs in several respects: it retains standard GR on the left-hand side (the geometry is not modified); it provides an explicit stress--energy interpretation; it naturally accommodates both compression and rarefaction phases; and it admits a cosmological completion route. MOND's relativistic extension TeVeS \citep{Bekenstein2004} introduces a scalar and vector field---one of our completion routes (\S3.5i) is structurally similar but derived from an inverse-GR starting point rather than a modified-gravity one.

\subsection{Relation to Emergent Gravity and Nonlocal Gravity}

Verlinde's emergent gravity program \citep{Verlinde2017} derives a MOND-like force from entropic considerations. The response-sector approach is compatible but more conservative: we do not assume a thermodynamic origin for gravity, only that the gravitational sector admits an effective constitutive response. Nonlocal gravity models \citep{DeserWoodard2007} can be rewritten as effective sources on the right-hand side of the Einstein equation, placing them within our bookkeeping framework.

\subsection{Strengths of the Current Work}

The primary contributions of this paper are:
\begin{itemize}
	\item a mathematical framework (inverse-GR $\to$ response-sector closure) that makes the ``alternative to CDM'' concept precise within standard GR;
	\item systematic testing against 175 galaxies spanning \textasciitilde{}$10^4$ in luminosity;
	\item independent confirmation of fitted amplitudes by a robust estimator;
	\item explicit falsification criteria stated in advance; and
	\item a domain-partitioned roadmap that acknowledges what is and is not claimed.
\end{itemize}

\subsection{Limitations and Open Questions}

The framework as presented has clear limitations. The boundary-layer model is phenomenological: the precise form of \textbf{(see PDF; boundary-layer functions not recovered)} and \textbf{(see PDF; boundary-layer functions not recovered)} is not derived from first principles. The lensing prediction is a working hypothesis, not a tested result. Cosmological viability is stated as a requirement, not demonstrated. Galaxy clusters---where the CDM paradigm is well-tested---are not addressed. Pressure-supported systems (elliptical galaxies, galaxy groups) require extension beyond the circular-orbit framework. These limitations define the research program going forward.

\subsection{The Historical Lesson}

The Vulcan episode provides a cautionary tale: gravitational anomalies initially attributed to unseen matter ultimately revealed deeper geometric structure. We do not claim the contemporary analogy is exact---the dark-matter hypothesis is vastly more successful and well-motivated than Vulcan ever was---but the structural lesson remains: when a gravitational anomaly is observed, both matter-based and geometry-based explanations deserve rigorous investigation. The persistent null results in dark-matter searches lend urgency to this investigation.

\subsection{A Candidate Identity for the Cold Dark Matter Role}

\subsubsection{Role vs. Identity: A Formal Distinction}

The Lambda Cold Dark Matter ($\Lambda$CDM) model stands as the standard paradigm of cosmology, a testament to its success in explaining phenomena from the Cosmic Microwave Background (CMB) to large-scale structure \citep{Planck2018VI2020}. Yet, as Kuhn argued, the accumulation of significant anomalies is the necessary precondition for scientific revolution \citep{Kuhn1962}. For $\Lambda$CDM, the central anomaly is existential: after decades of engagement, every attempt to directly identify the foundational particle of cold dark matter has returned a null result \citep{UndagoitiaRauch2016,LZ2023}.

This motivates a formal distinction between the role of dark matter and its identity. This distinction is critical for a productive scientific dialogue.

\subsection{Defining the CDM Role}

In the language of General Relativity (GR), the CDM role is the assertion that there exists an effective stress-energy tensor sector, \textbf{(see PDF; CDM-role sector symbol not recovered)}, such that the Einstein Field Equations are satisfied:

\textbf{(see PDF; Einstein-equation display equation not recovered from DOCX extraction)}

This combined tensor must yield a metric that reproduces the full suite of geodesic phenomenology: time-like geodesics (galaxy and cluster dynamics), null geodesics (gravitational lensing), and cosmological perturbations (structure growth and the CMB power spectrum) \citep{Planck2018VI2020}. The minimal constraints on this role are that the total stress-energy is conserved \textbf{(see PDF; conservation equation not recovered)} and that the resulting phenomenology is consistent across all observational domains.

\subsection{The Intrinsic Response as a Candidate Identity}

The intrinsic response framework proposes a specific identity to fill this role. It posits that the extra sector is not an independent entity with its own initial conditions (i.e., a particle fluid), but is instead a constitutive response functional of the baryonic sector itself:

\textbf{(see PDF; constitutive-response equation not recovered from DOCX extraction)}

This claim is profound: it suggests the gravitational phenomenon of dark matter is real, but it is an emergent, structured response of the spacetime metric to baryonic matter. The ``darkness'' is simply that this response is non-luminous and need not have electromagnetic couplings. This reframes the 50-year null detection record not as a failure to find dark matter, but as a stunning success in defining what it is not: a particle with non-gravitational interactions. Intrinsic response arguably elegantly fits both these ``nots'' as well as the phenomenology. The null results have carved out a residual space of possibility shaped exactly like a geometric, field-theoretic response---a puzzle piece for which the intrinsic response is a precise fit.

This is not a competitor to dark matter; it is a candidate for its identity. It avoids the primary vulnerability of Modified Newtonian Dynamics (MOND) and its relativistic extensions like TeVeS, which must explain away the evidence for a dark component on cluster and cosmological scales \citep{Milgrom1983,Bekenstein2004}. Instead, it seeks to be that component.

\subsubsection{A Falsifiable Validation Plan: The Discriminants Ladder}

A candidate identity warrants engagement only if it is falsifiable. The intrinsic response framework offers a clear validation plan in the form of a ``discriminants ladder,'' where each rung represents a decisive empirical test. A crucial first step is to declare a default theoretical branch for the galaxy regime:
\begin{itemize}
	\item \textbf{Branch A (Metric-Equivalent):} The scalar potentials are equal \textbf{(see PDF; condition not recovered)}, implying no gravitational slip. Lensing follows dynamics in the standard way.
	\item \textbf{Branch B (Slip-Allowed):} The potentials differ \textbf{(see PDF; condition not recovered)} in the boundary-layer regions, predicting a specific lensing-dynamics mismatch.
\end{itemize}

\subsubsection{Development Gateway: Branch Selection as an Explicit Research Fork}

The framework does not treat Branch A vs.
Branch B as a philosophical preference; it treats this as a development gateway that future work must adjudicate empirically.

If Branch A (metric-equivalent, \textbf{see PDF; no-slip condition}) holds in the galaxy regime: then lensing is a cross-check of the same metric potential inferred from dynamics, and the response sector behaves as a metric-equivalent ``dark gravity'' identity.

If Branch B (slip-allowed, \textbf{see PDF; slip condition} in boundary layers) is supported by data: then the response sector predicts a controlled, baryon-conditioned mismatch between dynamical and lensing inferences, providing a sharper discriminant and a tighter handle on covariant completion.

In both cases, failure occurs only if the required behavior cannot be achieved without reintroducing an independent, freely specified pressureless dark fluid (i.e., restoring the same degrees of freedom intrinsic response is meant to compress).

This ladder provides a clear, pre-registered path to validation or falsification, moving the debate from philosophical preference to empirical adjudication.

Adopting Branch A as the default, the ladder of falsifiers is as follows:

\textbf{(see PDF; Table 2 ``Rung Table (Branch A)'' not recovered from DOCX extraction)}

Branch B (Slip-Allowed): We allow a localized gravitational slip,
\textbf{(see PDF; slip functional form not recovered)}
in boundary/activation regions (e.g., near \textbf{see PDF; transition radius symbol}, within width \textbf{see PDF; width symbol}), while recovering \textbf{(see PDF; recovery condition)} in regimes where GR is tightly tested (Solar System) and/or where the response is inactive.

Adopting Branch B as the default, the ladder of falsifiers is as follows:

\textbf{(see PDF; Table 3 ``Rung Table (Branch B)'' not recovered from DOCX extraction)}

Branch B note: In this branch, weak and strong lensing are not expected to be identical to a metric-equivalent mapping from rotation curves; instead, lensing becomes the primary probe of slip structure and its activation rules.

\subsubsection{From Critique to Collaboration: A Playbook for Scientific Engagement}

The framework's incomplete status is not a vulnerability but a well-defined research opportunity, characteristic of a progressive research program in the sense of Lakatos \citep{Lakatos1978}. It invites collaboration by turning common criticisms into testable questions, as summarized in the following playbook:

\textbf{(see PDF; Table 4 ``Criticism and Rebuttal'' not recovered from DOCX extraction)}

\subsubsection{Conclusion: The Case for Engagement}

The intrinsic response framework offers a compelling synthesis in the long-standing dark matter debate. It reframes the problem by proposing a candidate identity---an emergent metric response---that is shaped by the very null results that have placed the particle hypothesis under strain. It is ontologically parsimonious, empirically grounded in the tight kinematic correlations of galaxies \citep{McGaughLelliSchombert2016}, and structured as a progressive and falsifiable research program.

The framework is incomplete, but it is not vague. It has moved beyond speculation to make specific, testable predictions that distinguish it from both standard $\Lambda$CDM and MOND-like alternatives \citep{BullockBoylanKolchin2017,Verlinde2017}. For these reasons, it has earned the right to serious scientific engagement. The question is no longer whether it is too incomplete to consider, but whether it is too compelling to ignore.


